
Current advances in information technology demand the integration of digital systems within educational institutions to improve the efficiency of academic administration. However, the implementation of Internships at SMK PGRI Telagasari still faces significant managerial challenges due to reliance on conventional procedures. The process of recording attendance, completing student logbooks, and summarizing grades still relies heavily on physical documents. This limitation creates various operational risks, such as reporting delays, data loss, and slow administrative verification, which makes it difficult for school management to conduct rapid evaluations crucial for strategic decision-making.

To address these issues, this final project develops InternSync, a management information system for monitoring Internships. This system is built using a centralized REST API-driven client-server architecture that integrates modern technologies, including React on the web frontend, the Laravel framework on the backend, and a MySQL database. A mobile application using the Flutter framework is still at the planning stage and is proposed as future development. InternSync provides Role-Based Access Control for five primary school roles alongside secure token-based access for industry mentors. The system integrates an internship vacancy catalog managed by coordinators and selected by students with coordinator screening, digital daily logbooks, attendance permission workflows with flexible industry confirmation and school fallback mechanisms, as well as final evaluations based on standardized vocational rubrics. The InternSync system was comprehensively tested using the Black Box Testing method, focusing on validating input and output functionality in a public deployment environment. The test results showed that all core functionality parameters, including CRUD operations in the submission module, industry partner master data management, logbook monitoring, error handling mechanisms, and interface responsiveness across screen sizes on the web application, ran successfully. The InternSync system has proven reliable in accelerating bureaucratic flows and optimizing the real-time monitoring function of PKL to support the implementation of the Smart School concept at SMK PGRI Telagasari.

\textbf{Keywords:} Management Information System, Internship Monitoring, Laravel, React, REST API.
